\documentclass[12pt,a4paper]{article}
\usepackage[T1]{fontenc}
\usepackage[utf8]{inputenc}
\usepackage{mathpazo}
\usepackage{amsmath,amssymb}
\usepackage[a4paper,margin=2.5cm]{geometry}
\usepackage{setspace}
\usepackage{booktabs,array}
\usepackage{xcolor}
\usepackage[most]{tcolorbox}
\usepackage{titlesec}
\usepackage{enumitem}
\usepackage[colorlinks=true,linkcolor=accent,urlcolor=accent]{hyperref}

\definecolor{accent}{HTML}{1F4E79}
\definecolor{accent2}{HTML}{8A4B08}

\setstretch{1.3}
\setlength{\parskip}{0.4em}
\setlength{\parindent}{0pt}

\titleformat{\section}{\Large\bfseries\color{accent}}{\thesection}{0.6em}{}
\titleformat{\subsection}{\large\bfseries\color{accent2}}{\thesubsection}{0.6em}{}
\titlespacing*{\section}{0pt}{2.2ex plus 0.5ex}{1ex}
\titlespacing*{\subsection}{0pt}{1.8ex plus 0.4ex}{0.6ex}

\newtcolorbox{keybox}[1][]{colback=accent!5,colframe=accent,boxrule=0.9pt,arc=2pt,
  left=8pt,right=8pt,top=5pt,bottom=5pt,#1}

\setlist[itemize]{itemsep=0.1em,topsep=0.2em}
\setlist[enumerate]{itemsep=0.1em,topsep=0.2em}

\newcommand{\E}{\mathbb{E}}
\newcommand{\tabfont}{\small\setstretch{1.05}}

\begin{document}

\begin{center}
{\Huge\bfseries\color{accent} EPL Match Outcome Predictor}\\[0.6em]
{\Large Checkpoint 1: Problem, Data Audit and Baselines}\\[0.6em]
{\large Dhruv Upadhyaya \quad$\cdot$\quad 20 October 2026 \quad$\cdot$\quad End of Phase 1}
\end{center}

\vspace{0.5em}

\begin{keybox}
\textbf{Where we are.} Twelve seasons of Premier League results and odds (4{,}240 matches) are
downloaded, cleaned and audited, and the data is sound. The bar to beat is a log loss of about
$0.95$ to $0.97$, which is what bookmaker odds achieve on 2023/24 to 2025/26. The plan
for Phase 2 is a from-scratch Poisson model with time decay, fitted by maximum likelihood with
hand-derived gradients.
\end{keybox}

%% ------------------------------------------------------------------
\section{Problem statement}

For every Premier League match $i$ we observe the date $t_i$, the home team $h(i)$, the away team
$a(i)$, the full-time goals $(x_i,y_i)$ and the result $r_i\in\{H,D,A\}$. The task is to output,
\emph{before kickoff}, a probability vector
\[
\mathbf p_i=(p_i^{H},\,p_i^{D},\,p_i^{A}),\qquad p_i^{k}\ge 0,\qquad p_i^{H}+p_i^{D}+p_i^{A}=1,
\]
using only matches $j$ with $t_j<t_i$. We predict probabilities, not just a winner, because
football is noisy: a hard prediction throws away how confident the model is, and probabilities
can be checked for calibration.

\subsection{How we score a model}

Let $N$ be the number of test matches. Three metrics are used.
\[
\mathrm{LL}=-\frac{1}{N}\sum_{i=1}^{N}\log p_i^{\,r_i}
\]
\[
\mathrm{BS}=\frac{1}{N}\sum_{i=1}^{N}\ \sum_{k\in\{H,D,A\}}\bigl(p_i^{k}-\mathbf 1[r_i=k]\bigr)^{2}
\]
\[
\mathrm{Acc}=\frac{1}{N}\sum_{i=1}^{N}\mathbf 1\bigl[\arg\max_k p_i^{k}=r_i\bigr]
\]
Log loss (LL) and the Brier score (BS) are strictly proper scoring rules: the expected score is
minimised only by reporting the true probabilities, so a model cannot gain by exaggerating its
confidence. Accuracy is reported for readability but is not used to fit or select models, since it
ignores probabilities entirely.

\begin{keybox}
\textbf{Reference values for a model with no information} (all three outcomes at $\tfrac13$):
\[
\mathrm{LL}=\log 3\approx 1.0986,\qquad
\mathrm{BS}=\Bigl(1-\frac13\Bigr)^{2}+2\Bigl(\frac13\Bigr)^{2}=\frac{6}{9}=\frac23 .
\]
\end{keybox}

%% ------------------------------------------------------------------
\section{Data and audit}

\textbf{Source.} football-data.co.uk, Premier League files \texttt{E0\_<season>.csv}, seasons
2015/16 to 2026/27. Full seasons have 380 matches each, so
\[
11\times 380 + 60 = 4{,}240 \text{ matches},
\]
where the 60 are the partial 2026/27 season (last match on 12 October 2026). The raw files are never
edited; \texttt{src/data.py} writes a cleaned table to
\texttt{data/\allowbreak processed/\allowbreak matches.csv}.

\textbf{Cleaning rules applied by the loader.}
\begin{itemize}
\item Files are read as \texttt{utf-8-sig}, because four of them start with a byte-order mark.
\item Dates are parsed day-first, accepting both \texttt{dd/mm/yy} and \texttt{dd/mm/yyyy}.
\item The \texttt{Time} column is not required, since it is absent in 2015/16 to 2018/19.
\item Only core columns and the B365, Average and Pinnacle (PS) odds are kept. Same-match
statistics such as shots and possession are dropped on purpose: they are unknown before kickoff and
would leak the result into the features.
\end{itemize}

\textbf{Structural checks, all passed.}
\begin{itemize}
\item The result column agrees with the goals in every match, and there are no duplicate fixtures.
\item Every full season has 20 teams and 19 home games per team.
\item Every season transition has exactly 3 promoted and 3 relegated teams.
\item Team names are spelled consistently across seasons (35 distinct names in total). The only
similar pair flagged, West Brom and West Ham, is a false positive.
\end{itemize}

\subsection{Findings that affect the modelling}

\begingroup
\tabfont
\renewcommand{\arraystretch}{1.25}
\begin{tabular}{@{}>{\raggedright\arraybackslash}p{0.43\linewidth}>{\raggedright\arraybackslash}p{0.5\linewidth}@{}}
\toprule
\textbf{Finding} & \textbf{Consequence}\\
\midrule
Average odds are missing in all 1{,}520 matches of 2015/16 to 2018/19. &
Odds baselines can only be computed from 2019/20 onward with Average odds; B365 is complete.\\
PS odds stop on 8 January 2026. The 170 matches from 17 January to 24 May 2026, and all of 2026/27,
have none. &
PS comparisons use only matches that have PS odds (970 in the test window).\\
2026/27 is partial: 60 matches, 3 home games per team. &
Excluded from the baselines. It supplies recent form only.\\
2020/21: away wins 40.3\% against 37.9\% home wins. &
Supports giving closed-doors matches their own home-advantage term.\\
Mean goals jump in 2023/24 to 1.80 (home) and 1.48 (away), against about 1.5 and 1.2 earlier. &
The league-wide goal rate is not constant, so time decay or a per-season intercept is needed.\\
\bottomrule
\end{tabular}
\endgroup

%% ------------------------------------------------------------------
\section{Baselines}

\subsection{Bookmaker probabilities}

Decimal odds $o_H,o_D,o_A$ imply probabilities $1/o_k$ that sum to more than one. The excess is the
bookmaker's margin (the overround),
\[
\Omega=\sum_{k\in\{H,D,A\}}\frac{1}{o_k}-1 .
\]
Removing it by proportional normalisation gives
\[
p^{k}=\frac{1/o_k}{\displaystyle\sum_{j\in\{H,D,A\}}1/o_j}.
\]
The mean margin is 5.5\% for B365, 4.8\% for Average and 3.5\% for PS. Proportional normalisation
is the simplest choice; it ignores the favourite--longshot bias, and alternatives (Shin's method, the
power method) can be tried later.

\subsection{Results on 2023/24 to 2025/26}

All baselines are scored on the last three completed seasons (1{,}140 matches). The class
frequencies for the second row come from 2015/16 to 2022/23 only, so no test information is used.

\begin{center}
\tabfont
\renewcommand{\arraystretch}{1.25}
\begin{tabular}{lrrr}
\toprule
\textbf{Predictor} & \textbf{Matches} & \textbf{Log loss} & \textbf{Accuracy}\\
\midrule
Uniform, $\tfrac13$ each & 1{,}140 & 1.0986 & 43.2\%\\
Training class frequencies (H 0.448, D 0.234, A 0.318) & 1{,}140 & 1.0728 & 43.2\%\\
B365, margin removed & 1{,}140 & 0.9662 & 54.2\%\\
Average, margin removed & 1{,}140 & 0.9650 & 54.2\%\\
Pinnacle (PS), margin removed & 970 & 0.9501 & 55.6\%\\
\bottomrule
\end{tabular}
\end{center}

The accuracy in the first two rows equals that of always predicting a home win (43.2\%); that rule
has infinite log loss because it assigns probability zero to draws and away wins.

\begin{keybox}
\textbf{The bar to beat.} On the 970 matches that have all three sources, the bookmakers agree to
within $0.001$ (log loss $0.9501$ to $0.9510$). That subset is the fairest benchmark, because the PS
figure above is flattered by which matches it happens to cover. On the full 1{,}140 matches, the
Average-odds log loss of $0.965$ is the reference. Bookmakers also use injuries and team news, which
our model will not have, so matching them within about $0.01$ from past results alone would
already be a strong outcome.
\end{keybox}

%% ------------------------------------------------------------------
\section{Model roadmap and the mathematics}

\subsection{Model A: multinomial logistic regression (baseline)}

With a feature vector $\mathbf z_i$ built only from earlier matches (rolling form, an Elo
difference, home flag),
\[
p_i^{k}=\frac{\exp(\mathbf w_k^{\top}\mathbf z_i)}{\displaystyle\sum_{j\in\{H,D,A\}}\exp(\mathbf w_j^{\top}\mathbf z_i)} .
\]
Softmax is unchanged when the same vector is added to every $\mathbf w_k$, so one class is fixed as a
reference, $\mathbf w_A=\mathbf 0$. Minimising the cross-entropy (the log loss) gives the gradient
\[
\frac{\partial\,\mathrm{LL}}{\partial\mathbf w_k}=\frac{1}{N}\sum_{i=1}^{N}\bigl(p_i^{k}-\mathbf 1[r_i=k]\bigr)\mathbf z_i .
\]

\subsection{Model B: independent Poisson goals}

Model the goals of each match as independent Poisson variables, $X_i\sim\mathrm{Pois}(\lambda_i)$ for
the home side and $Y_i\sim\mathrm{Pois}(\mu_i)$ for the away side, with
\[
\log\lambda_i=\iota+\eta+\alpha_{h(i)}-\beta_{a(i)}
\]
\[
\log\mu_i=\iota+\alpha_{a(i)}-\beta_{h(i)} .
\]
Here $\iota$ is the league-wide log goal rate, $\eta$ is the home advantage, $\alpha_t$ is the attack
strength of team $t$ and $\beta_t$ its defensive strength (larger means fewer goals conceded). The
model has $2+2T$ parameters for $T$ teams.

\textbf{Identifiability.} The rates $\lambda_i$ and $\mu_i$ depend on $\iota$ and on the averages
$\bar\alpha$ and $\bar\beta$ of the strengths only through the combination $\iota+\bar\alpha-\bar\beta$.
So the likelihood alone cannot separate two directions of change: adding $c$ to every $\alpha_t$ and
every $\beta_t$, and adding $d$ to $\iota$ while subtracting $d$ from every $\alpha_t$. A ridge penalty
on $\alpha$ and $\beta$ (but not on $\iota$) removes both. Every match has exactly one home and one
away team, so summing the gradient equations given below over teams gives
\[
\sum_t\frac{\partial\ell}{\partial\alpha_t}=\frac{\partial\ell}{\partial\iota}-2\kappa\sum_t\alpha_t
\]
\[
\sum_t\frac{\partial\ell}{\partial\beta_t}=-\frac{\partial\ell}{\partial\iota}-2\kappa\sum_t\beta_t .
\]
At the optimum every gradient vanishes, so $\sum_t\alpha_t=\sum_t\beta_t=0$ and $\iota$ carries the whole
league-wide level. The same penalty is a Gaussian prior that shrinks teams with little data (such as
newly promoted ones) towards the league average.

\textbf{Weighted penalised log-likelihood.} Let $T^{*}$ be the prediction date, measured in years, and
let
\[
w_i=e^{-\xi\,(T^{*}-t_i)},\qquad \tau_{1/2}=\frac{\ln 2}{\xi},
\]
so recent matches count more and $\tau_{1/2}$ is the half-life of a match's influence. With
$\theta=(\iota,\eta,\boldsymbol\alpha,\boldsymbol\beta)$,
\[
\ell(\theta)=\sum_i w_i\Bigl[x_i\log\lambda_i-\lambda_i+y_i\log\mu_i-\mu_i-\log(x_i!\,y_i!)\Bigr]
-\kappa\sum_t\bigl(\alpha_t^2+\beta_t^2\bigr).
\]

\textbf{Gradient.} Because $\log\lambda_i$ and $\log\mu_i$ are linear in $\theta$, we have
$\partial\lambda_i/\partial\theta_k=\lambda_i\,\partial\log\lambda_i/\partial\theta_k$, and the
Poisson terms collapse to residuals. Define the home and away residuals
\[
e_i=x_i-\lambda_i,\qquad f_i=y_i-\mu_i .
\]

\begin{keybox}
\vspace{-0.6em}
\begin{align*}
\frac{\partial\ell}{\partial\iota}&=\sum_i w_i\,(e_i+f_i)\\[2pt]
\frac{\partial\ell}{\partial\eta}&=\sum_i w_i\,e_i\\[2pt]
\frac{\partial\ell}{\partial\alpha_t}&=\sum_{i:\,h(i)=t}w_i\,e_i+\sum_{i:\,a(i)=t}w_i\,f_i-2\kappa\,\alpha_t\\[2pt]
\frac{\partial\ell}{\partial\beta_t}&=-\sum_{i:\,a(i)=t}w_i\,e_i-\sum_{i:\,h(i)=t}w_i\,f_i-2\kappa\,\beta_t
\end{align*}
\vspace{-1.2em}
\end{keybox}

Setting these to zero has a clear reading: without the penalty, the weighted totals of expected
goals must equal the weighted totals of observed goals, overall (for $\iota$), at home (for $\eta$),
and for each team's attack and defence.

\textbf{Concavity.} Write $\mathbf a_i$ and $\mathbf b_i$ for the (sparse, $0/\pm1$) coefficient vectors
of $\log\lambda_i$ and $\log\mu_i$ in $\theta$. The Hessian is
\[
\nabla^2\ell=-\sum_i w_i\bigl(\lambda_i\,\mathbf a_i\mathbf a_i^{\top}+\mu_i\,\mathbf b_i\mathbf b_i^{\top}\bigr)-2\kappa\,I_{\alpha,\beta},
\]
which is negative semi-definite. The two flat directions of the likelihood described above both change
$\alpha$ or $\beta$, where the ridge term is strictly negative, so $\ell$ is strictly concave. It
therefore has a unique maximiser, and a quasi-Newton method such as L-BFGS reaches it. The analytic
gradient is checked against finite differences in
\texttt{tests/\allowbreak test\_poisson.py}.

\textbf{From goals to outcomes.} For independent Poisson goals,
\[
P(X=x,Y=y)=\frac{e^{-\lambda}\lambda^{x}}{x!}\cdot\frac{e^{-\mu}\mu^{y}}{y!},
\]
and the three outcome probabilities are sums over the score grid,
\[
p^{H}=\sum_{x>y}P(x,y),\qquad p^{D}=\sum_{x=y}P(x,y),\qquad p^{A}=\sum_{x<y}P(x,y),
\]
truncated at about ten goals per side, where the lost probability is negligible. For the independent
model, $X-Y$ follows a Skellam distribution, which gives an exact cross-check of the truncated sums.

\subsection{Model C: Dixon--Coles correction}

Independent Poisson models misprice low scores, especially draws. Dixon and Coles multiply the
probability of the four lowest scorelines by a factor $\tau$:
\[
\tau(x,y)=
\begin{cases}
1-\lambda\mu\rho & (x,y)=(0,0)\\
1+\lambda\rho & (x,y)=(0,1)\\
1+\mu\rho & (x,y)=(1,0)\\
1-\rho & (x,y)=(1,1)\\
1 & \text{otherwise.}
\end{cases}
\]
All factors must stay positive, which restricts the dependence parameter to
\[
\max\!\Bigl(-\frac{1}{\lambda},-\frac{1}{\mu}\Bigr)\le\rho\le\min\!\Bigl(\frac{1}{\lambda\mu},\,1\Bigr).
\]
The correction conserves total probability. Using $P(0,0)=e^{-\lambda-\mu}$,
$P(0,1)=e^{-\lambda-\mu}\mu$, $P(1,0)=e^{-\lambda-\mu}\lambda$ and $P(1,1)=e^{-\lambda-\mu}\lambda\mu$,
the net change is
\[
e^{-\lambda-\mu}\,\rho\,\bigl(-\lambda\mu+\lambda\mu+\lambda\mu-\lambda\mu\bigr)=0 .
\]
The log-likelihood gains the term $\log\tau(x_i,y_i)$ for each match, and $\rho$ is fitted together
with $\theta$. The time-decay rate $\xi$ and the ridge strength $\kappa$ are chosen by validation,
not set by hand.

\subsection{Evaluation protocol}

Models are compared by \emph{walk-forward validation}: for each block of matchweeks in the test
window, fit on all earlier matches, predict the block, then move forward. This mimics real use,
avoids leakage, and gives about 1{,}140 out-of-sample predictions for the 2023/24 to 2025/26 window.
Random shuffling of matches is never used.

%% ------------------------------------------------------------------
\section{Open decisions and risks}

\begin{enumerate}
\item \textbf{Which odds set is the benchmark.} The columns kept so far are pre-match odds. Recent
football-data.co.uk files also carry closing-odds columns, which are usually sharper. If they are
present in our files, they define a stricter bar, and the loader should be extended.
\item \textbf{Closed-doors window.} The home-advantage indicator should be defined by match date,
covering the restart of 2019/20 in June 2020 and most of 2020/21, with exceptions where limited
crowds were allowed. The exact dates must be confirmed before the feature is built.
\item \textbf{The 2023/24 goal-rate shift.} Options are time decay alone, or an extra per-season
intercept. Validation decides.
\item \textbf{Team-specific home advantage.} It adds $T$ more noisy parameters. We will test it
against a single $\eta$ rather than assume it helps.
\item \textbf{Draws are hard.} Even good models rarely make draws the most likely outcome, so
accuracy will understate model quality. Log loss and Brier score are the primary metrics.
\end{enumerate}

%% ------------------------------------------------------------------
\section{Phase 2 plan}

\begin{enumerate}
\item Write the Poisson model in \texttt{src/poisson.py}: log-likelihood and analytic gradient,
with a finite-difference gradient check.
\item Simulate matches from a Poisson model with known parameters, fit, and confirm the parameters are
recovered and the error shrinks as the sample grows.
\item Build Model A with rolling-form and Elo features, as the first real comparison.
\item Fit Model B on the real data and run walk-forward validation against the baselines above.
\item Add the Dixon--Coles term and time decay (Model C); choose $\xi$ and $\kappa$ by validation.
\item Keep a log of every bug and fix for the final report.
\end{enumerate}

\vspace{0.5em}
\begin{center}
\small\color{gray} Figures in this report are taken from the output of \texttt{src/audit.py}
(\texttt{reports/checkpoint1\_tables.md}).
\end{center}

\end{document}
